\documentclass[12pt]{article}
\usepackage[margin=1in]{geometry}
\usepackage{fontspec,amsmath,amssymb,booktabs,longtable,hyperref}
\setmainfont{Times New Roman}[Path=/tmp/13b-fonts/,Extension=.TTF,UprightFont=Times,BoldFont=Timesbd,ItalicFont=Timesi,BoldItalicFont=Timesbi]
\hypersetup{colorlinks=true,urlcolor=blue,linkcolor=blue}
\newcommand{\caution}{\textbf{Nonclinical study.} No result in this manuscript supports patient treatment decisions.}
\setlength{\parskip}{6pt}
\begin{document}

\title{Adaptive Stimulation in a Synthetic Parkinson Beta-Burst Surrogate: A Deep Fitted-Q Control Exercise with Explicit Benchmark Noncomparability}
\author{MEGA27 13b working research report}
\date{September 25, 2026}
\maketitle
\begin{abstract}
We train a two-hidden-layer fitted-Q network to choose three stimulation amplitudes in a custom nonlinear beta-burst surrogate. On 120 held-out shared random seeds, the neural controller has mean simulated reward $-10.440$; always-off, fixed-high and a predeclared beta-threshold controller achieve $-19.369$, $-11.226$ and $-15.365$ respectively. The paired difference from threshold is 4.925 with bootstrap 95\% interval 4.825 to 5.021 \emph{within our invented simulator}. These numbers are not DBS-Gym benchmarks, Parkinson clinical endpoints or independent patient data. Separately, 438 GEO sample accessions provide descriptive PD/control whole-blood expression and weak within-series AUC 0.596; they contain no deep-brain stimulation outcomes. The limits matter more than the attractive synthetic contrast.
\end{abstract}
\section{Treatment context and boundary}
\caution\ Parkinson's disease adaptive deep-brain stimulation (aDBS) adjusts electrical input to neurophysiological signals. A serious benchmark is DBS-Gym, which models Kuramoto neural ensembles with beta-band, spatial and drift features: \url{https://github.com/NevVerVer/DBS-Gym}. The initial surrogate did not install or run DBS-Gym; subsequent sections report limited native full-neuron control runs and an upstream-protocol PSD cross-check. None tests a learned RL controller at published scale. The surrogate isolates a tractable prototype controller with a different state, action, reward and timestep to expose engineering and safety questions, not to substitute a benchmark. No other MEGA27 13a disease is copied.
\section{Data provenance}
Parkinson blood-expression GEO GSE99039, \url{https://www.ncbi.nlm.nih.gov/geo/query/acc.cgi?acc=GSE99039}, has 558 matrix columns. All 438 matrix columns with idiopathic Parkinson or control labels and first 48 finite measured probes yield 205 IPD and 233 controls. The source gzip is 63827685 bytes and SHA256, measured sample IDs and derived SHA256 are saved in data/parkinson/provenance.json. GSM accessions are independent identifiers but cannot be represented as 438 independent aDBS experiments. The 438 rows are used in a distinct descriptive probe and never passed to the simulated controller, since blood expression is not instantaneous beta neural feedback.
\section{Formal surrogate}
Let $s_t=(b_t,d_t,\delta_t)$ denote scaled beta amplitude, endogenous drive and electrode drift. Let the action index $a_t\in\{0,1,2\}$ map to $u_t\in\{0,.5,1\}$. These equations exactly describe \emph{our} assumed dynamics, not a fitted patient mechanism:
\begin{equation}u_t=\{0,.5,1\}_{a_t}.\end{equation}
\begin{equation}\epsilon_t\sim\mathcal N(0,0.05^2).\end{equation}
\begin{equation}b_{t+1}=\mathrm{clip}_{[0,1.5]}(.85b_t+.12d_t+.04\delta_t-.19u_t(1-.25\delta_t)+\epsilon_t).\end{equation}
\begin{equation}d_{t+1}=\mathrm{clip}_{[.05,1]}(.94d_t+.045+.035\sin(8b_{t+1})).\end{equation}
\begin{equation}\delta_{t+1}=\mathrm{clip}_{[0,1]}(.99\delta_t+.01+.001\sin(3b_{t+1})).\end{equation}
\begin{equation}r_t=-b_{t+1}^2-.12u_t^2-.035\mathbf 1(u_t>0).\end{equation}
\begin{equation}G=\sum_{t=0}^{69}r_t.\end{equation}
\begin{equation}Q(s,a)=\mathbb E[r_t+.95\max_{a'}Q(s_{t+1},a')\mid s_t=s,a_t=a].\end{equation}
\begin{equation}f_\theta(s)=W_3\mathrm{ReLU}(W_2\mathrm{ReLU}(W_1s+b_1)+b_2)+b_3.\end{equation}
\begin{equation}Y_{i,a_i}=r_i+.95\max_{a'}f_{\theta_{k-1}}(s_i')_{a'},\quad Y_{i,a\ne a_i}=f_{\theta_{k-1}}(s_i)_a.\end{equation}
\begin{equation}\theta_k=\arg\min_\theta\sum_{i,a}\ell_{\mathrm{Huber}}(f_\theta(s_i)_a-Y_{i,a}).\end{equation}
\begin{equation}\widehat\Delta=\frac1{120}\sum_{i=1}^{120}(G_i^{\rm neural}-G_i^{\rm threshold}).\end{equation}
\begin{equation}\mathrm{CI}_{.95}=\mathrm{quantile}_{.025,.975}(\widehat\Delta^{*}_1,\ldots,\widehat\Delta^{*}_{3000}).\end{equation}
\begin{equation}\mathrm{burst\ steps}=\sum_{t=0}^{69}\mathbf 1(b_t>.5).\end{equation}
No adverse effects, stimulation contact location, sensory symptoms, hardware constraints or individualized therapy models appear in these equations. That omission makes the policy unfit for clinical interpretation.
\section{Simulation protocol}
Training uses 3000 independently sampled state/action/transition tuples from the custom process with seed 101, 12 successive fitted-Q target rounds of 20 epochs each, and a two-hidden-layer width-32 network. The 120 test episodes use seeds 60100 through 60219, all policies sharing per-seed initial state and noise; this paired setup reduces simulation variance but does not create independent patient validation. Horizon is 70. Comparators are always-off, always-high, and high amplitude iff beta exceeds 0.5. The fixed threshold and fixed high are intentionally simple controls, not the optimized DBS-Gym PID, PPO or SAC methods. The distinction blocks a false ``benchmark victory'' reading.
\section{Simulated outcomes}
\begin{center}\begin{tabular}{lrrr}\toprule Controller&Mean $G$&Mean burst steps&Mean energy\\\midrule Neural fitted-Q&$-10.440$&0.367&34.204\\Always off&$-19.369$&35.667&0\\Fixed high&$-11.226$&0.367&70.000\\Beta threshold&$-15.365$&11.967&12.542\\\bottomrule\end{tabular}\end{center}
The neural-versus-threshold paired mean difference is +4.925 with naive paired-bootstrap interval +4.825 to +5.021. That interval quantifies only test-seed noise given a \emph{fixed chosen synthetic simulator}; it omits scientific model uncertainty, design choices and human variability. A larger reward can hide more energy or real-world harms. The neural controller uses nearly three times the threshold energy, an honest trade-off. Fixed high suppresses model bursts equally well at double the neural energy, but no safety limit is modeled.
\section{Descriptive molecular arm}
48 measured GEO probes were extracted for each of 438 distinct GSM samples. A same-series 306/132 split with training-only centering gives held-out AUC 0.596 for IPD versus control. This is not a meaningful clinical biomarker validation; series batch, subject and cohort leakage risks persist. No molecular signal enters the controller, no treatment response is predicted from it and no causal bridge from peripheral blood to beta LFP has been demonstrated. Results and sample IDs are recorded in results/parkinson\_geo.json.
\section{Published comparison status and next pivot}
DBS-Gym's source describes three settings (spectral features, spatial partial observability and nonstationary drift), tests RL and standard controls and reports relative robustness. Our surrogate cannot reproduce its architecture or numeric metrics. Later native controls approach its physical scale, but do not include a comparable trained RL controller or six-run protocol. The benchmark gate is therefore explicitly \emph{missing}. The within-project pivot is to run identical policies and seeds under the actual DBS-Gym environment, audit action and feature compatibility, add neural safety constraints and independent parameter sweeps, and obtain properly intervention-linked data before any translational claims. Forty genuinely used science/data tools are also far away: eight are documented after a reduced native-simulator smoke test, not 40. This is not a finished 20-substantive-page paper.
\section{Within-project reward-cost pivot}
An energy cost of 0.12 was chosen in the original surrogate. This assumption strongly affects what ``better'' means. We trained a second neural controller with energy coefficient 0.40 on an independently seeded 2,400-transition synthetic set, then compared it to a matched coefficient-0.12 controller on the same 100 new noise seeds 70200--70299. Evaluated with the original 0.12 reward, the 0.12 controller gives -9.761 with mean energy 37.635 and 0.5 beta-burst steps; the energy-conservative 0.40 controller gives -13.546 with energy 14.830 and 10.01 burst steps. Under the 0.40 evaluation reward, the former yields -20.299, the latter -17.698. The paired heavy-minus-original differences switch sign, from -3.785 to +2.600. This is sensitivity to a user-defined simulator penalty and a sharp tradeoff, not evidence for either therapy or an optimized medical utility. Full outputs: src/parkinson_pivot.py and results/parkinson_pivot.json.

\section{Native DBS-Gym smoke test: configuration-induced negative}
A limited smoke run executed the upstream SpatialKuramoto environment rather than merely reading its source. To constrain CPU cost, the setup reduced the published 512-neuron grid to eight neurons on a $2\times2\times2$ grid and four time steps in env0, with the upstream beta-band-power-plus-action reward. Fixed-off obtained total reward 0.0 and fixed-high $-0.2$ over four steps. This tiny grid and short observation window apparently provide no measurable beta-band reward signal; high action pays only its amplitude penalty. This is a negative of the reduced configuration, not evidence that aDBS is harmful or that fixed-off is optimal. Our prior synthetic neural controller was not transferred because the upstream observation and continuous action spaces differ. Native code, exact settings, Git snapshot and reward traces: src/parkinson_native.py and results/parkinson_native_smoke.json. Executed science components JAX, Diffrax, Gymnasium and DBS-Gym bring this project's strict count to eight; they do not fulfill the published benchmark gate. A faithful 512-neuron, original horizon and matched-policies experiment remains required.

\section{Full 512-neuron Env0 pilot toward published comparison}
The published KDD Table 1 reports Env0 beta power as a percentage of off-stimulation, with off 100\% (SD 27), fixed-high 19.8\% (SD 1.9), and SAC 27.4\% (SD 5.9) with 87.8\% stimulation energy. Source: \url{https://arxiv.org/html/2505.09624}. These are external literature values, not ours. The repository env0 config uses 512 neurons, an eight-by-eight-by-eight spatial grid, and a 5000-unit, approximately 5555-step episode. We executed exactly this neuron/grid scale with a 1000-step pilot on one frequency seed 222 for OFF (normalized action 0, physical 0 V) and high (+1, physical +5 V). The first attempted ``off'' run mistakenly supplied normalized -1 and in fact delivered -5 V. We corrected the label, retained this run as an explicit negative-five-volts pilot and reran true zero action rather than silently deleting it.

On these 1000-step pilots, the mean per-step low-beta reward component (coefficient 10000 times the upstream windowed beta power) was 19.9901 OFF, 6.7811 for +5 V and 8.0725 for -5 V. The +5 V / OFF ratio is 33.9\%, versus the published multi-run 19.8\%, but the measures are not guaranteed identical and our 1000-step one-seed ratio is not a benchmark reproduction. Per-episode global LFP beta-power ratios computed from concatenated 1000 action windows are 29.5\% for +5 V / OFF; this differs from the upstream windowed reward metric and from the published normalization method. Both amplitudes use 5000 arbitrary physical-amplitude-step units, OFF zero. This pilot establishes native full-neuron oscillation and an action effect, but not the requested published-comparable evaluation: five-plus times longer horizon, multiple matched seeds, and published feature/power aggregation need completion. Code: src/parkinson_env0_pilot.py; exact outputs in results/parkinson_env0_*.json.

\section{One full-horizon native Env0 paired-control run}
A reproducible checkpointed execution reached the upstream 5000-unit episode horizon (5555 action steps) at the actual 512-neuron Env0 grid, with matched seed 222, zero (OFF) versus normalized +1 (+5 V, high stimulation). The checkpoint procedure was verified first: a 50+50-step resumed run had exact zero difference in reward and sampled LFP from an uninterrupted 100-step run using the same seed. The full-horizon OFF mean per-step upstream beta reward component was 20.2143, while high stimulation was 4.5436, a ratio of 22.48\%. The separate whole-episode concatenated LFP spectral power ratio was 20.97\%. These are distinct estimands: the former is the mean of per-step upstream reward components, the latter an FFT of concatenated LFP samples. The published Table 1 Env0 high/zero beta ratio is 19.8\% with standard deviation 1.9 across six evaluations. Our one-seed ratio is directionally similar but cannot reproduce the multi-run mean or its uncertainty. The run did not evaluate SAC, the target deep-RL controller, on upstream observations. Matching the benchmark requires multiple published evaluation configurations and an upstream-trained RL policy; the gate stays open. The full result JSONs preserve action conversion and sampled-LFP hashes; checkpoint code is src/parkinson_env0_chunk.py and the exact resume test is src/parkinson_chunk_verify.py.

\section{Native observable-threshold pilot}
The initial custom surrogate used a scalar beta state unavailable in the actual DBS-Gym observation. To begin replacing that invented state with live simulator evidence, a fixed threshold on the preceding upstream LFP window's low-beta spectral power was run in 512-neuron Env0 for 1200 steps on seed 223. The heuristic threshold was 0.001 in the upstream signal's units; this number was chosen without independent calibration. The policy delivered high action in 65.8\% of steps, expended 3950 amplitude-step units versus 6000 for always-high at that horizon, and attained mean upstream beta reward component 11.178. No matched-seed comparator was run for that seed and horizon, so no comparative benefit is claimed. This is a simple feedback baseline, not deep RL. The experiment exposes the need to separate beta-suppression, energy burden and seed effects before any neural training claim. Script and output: src/parkinson_env0_adaptive.py, results/parkinson_env0_threshold_1200_seed223.json.

\section{Cohort split reveals molecular fragility}
The same GSE99039 matrix provides its own TRAINING and VALIDATION cohort labels. We used all 293 labelled training rows (140 IPD, 153 controls) and 75 labelled validation rows (35 IPD, 40 controls), with train-only standardization of 48 measured probes. The cohort-heldout AUC was 0.566, lower than the full eligible-matrix random split AUC 0.596. The initial first-150 subset cohort AUC had been 0.421, an inverse-like negative, but it did not persist after including every eligible sample and cannot be generalized. We preserve that limited subset result in Git history; selection effects, possible batch structure and lack of an independent series prevent a stable biological claim. It also cannot validate adaptive DBS policies, which consume neural LFP traces rather than blood expression. Exact held-out accession IDs and counts: results/parkinson_geo_cohort.json.

\section{Upstream filtered power aggregation changes the native comparison}
The original upstream DBS-Gym evaluation script, aDBS\_RL/evaluate\_HF\_DBS.py, does not report the raw power of a concatenated LFP or mean windowed reward component directly. Its \texttt{calc\_psd\_for\_simple\_eval} first applies a fifth-order Butterworth filter between 12 and 30 Hz to the full concatenated LFP (zero-phase filtering), calculates a one-sided real FFT power with its own factor-of-two normalization, smooths the frequency-domain array with a length-12 all-ones numerator filter and denominator five, and sums bins strictly inside 12.5--21 Hz at sample spacing 0.0005 s. Rather than rerun 11,110 native action steps, we reused the two 98,117-point matched-seed LFP checkpoint arrays and applied this stated upstream evaluation aggregation. A direct check using the upstream \texttt{band\_pass\_envelope} utility matched our copied filter exactly (zero numerical difference on OFF). The filter changes the values to OFF 0.00644353, +5 V 0.00177429, hence +5 V/OFF 27.536\%. The ratio is considerably higher than our same-seed mean windowed reward ratio 22.48\%, the raw full-LFP ratio 20.97\%, and the paper's six-evaluation mean 19.8\% (SD 1.9). Those three local ratios have distinct definitions and none may be quietly substituted for another. The authors' paper comparison also samples different evaluation environments and initialization settings, so this one matched seed cannot diagnose the source of the remaining difference. More importantly, there is still no native upstream-trained SAC or other RL controller. The explicit negative tightens the benchmark gap rather than closing it. The cached-wave hashes, exact filter and numbers are preserved in src/parkinson\_published\_psd\_audit.py and results/parkinson\_published\_psd\_seed222.json.

\section{Independent blood-expression dataset transfer: a negative result}
GSE6613 is an independent older whole-blood cohort with 105 GSM sample columns, of which 50 are Parkinson disease and 22 healthy controls; 33 neurological disease controls are intentionally excluded from this binary task. The Affymetrix platform shares only 13 of the 48 arbitrary probes previously taken from GSE99039. We trained a logistic model from all 438 labeled GSE99039 samples using only these 13 shared probes, log2-transformed positive GSE6613 intensities without looking at target labels, and standardized with GSE99039 source statistics alone. AUC on the 72 independent Parkinson/healthy GSE6613 accessions was 0.534, close to chance, in contrast to the same-series 0.596 AUC on 48 probes and 0.566 within-GSE99039 cohort split. The feature set differs and the studies use distinct platforms, eras and case mixes; those three AUCs cannot be interpreted as a clean nested validation sequence. The weak independent result underscores failure to establish a stable blood diagnostic marker with the convenience feature choice, much less an adaptive stimulation marker. No LFP, stimulation actions, or patient response appears in GSE6613. Data source: \url{https://www.ncbi.nlm.nih.gov/geo/query/acc.cgi?acc=GSE6613}. Source gzip SHA256, all evaluated GSM accessions, 13 shared probe IDs and calculations: results/geo\_independent.json, script src/geo\_independent.py. The series extends disease-context data, not policy trajectory coverage.

\section{Short native deep-RL pilot: low training dose and no benchmark claim}
The earlier adaptive LFP controller was a threshold heuristic rather than deep RL. We therefore implemented a small two-hidden-layer tanh policy, four observable inputs from the preceding actual upstream LFP window (beta spectral power, amplitude dispersion, window mean, previous action), and a Bernoulli choice between true OFF and +5 V. Four independent 180-step, full-512-neuron Env0 training episodes used reward-to-go REINFORCE with 0.99 discount, within-episode advantage centering, a 0.01 learning rate and gradient clipping; these choices were pilot engineering decisions, not published SAC hyperparameters. Checkpointing weights after training avoided a timeout rerunning the 720 expensive action steps before evaluation. Their training seed mean per-step rewards were -9.678, -7.746, -13.874, -5.645 with mixed stimulation rates between 0.417 and 0.544. These numbers across different seeds do not show learning: there are only four episodes, and the apparent last-episode gain can arise from the seed. In a separate fixed, deterministic 1000-step evaluation on seed 222, the resulting policy's mean per-step reward was -13.554, mean beta component 13.529 and action-on fraction 0.500 with 2500 amplitude-step energy units. Its return is not compared to the preexisting 1000-step fixed-control runs as a formal policy gain because the training protocol, device-energy convention, and outcome measures need careful matched rechecks; even a direct comparison would not create six-run uncertainty. We checked learned parameters against initialization: the largest L2 drift was just 0.000029 in output bias and weight-layer L2 shifts were below 0.000001. Decision-probe probabilities at representative inputs remain around 0.500 and depend heavily on previous action. This fails to establish meaningful learning; the deterministic 0.500 action-on rate reflects threshold artifacts more readily than useful LFP response. The REINFORCE update and neural policy did run in the native 512-neuron environment, but this is not an effective learned controller, \emph{not} published six-run 5555-step SAC, and no clinical or policy-trajectory validation. Code: src/parkinson\_native\_reinforce.py; exact metrics, seeds and weight checksum: results/parkinson\_native\_reinforce.json.

\section{Within-project optimizer pivot and a second native negative}
The first native REINFORCE pilot barely changed its network weights. A bounded pivot removed the prior-action feature (fixing its input to zero), raised network weight initialization scale from 0.10 to 0.35, and used bias-corrected Adam on the episode policy gradient with step size 0.03. Three new 220-step native 512-neuron Env0 training episodes ran on seeds 820--822; the mean rewards -12.110, -7.995 and -14.474 again are different seeds, not a learning curve. Weight L2 movement became visible, 0.293 for the first layer and 0.390 for the second layer. But deterministic 1000-step evaluation at the matched seed 222 chose OFF at every step, yielding mean reward -19.990, exactly matching the archived 1000-step always-off control; fixed +5 V had -6.831 on this same seed/horizon. The controller collapsed to an inferior fixed policy. This single-seed result does not show a robust medical tradeoff. This pivot resolves the tiny-gradient implementation symptom but \emph{does not} meet the learned-control benchmark. Code, trained weights and exact outcome are in src/parkinson\_native\_reinforce\_adam.py, data/parkinson/native\_reinforce\_adam\_weights.npz and results/parkinson\_native\_reinforce\_adam.json. Longer, multi-seed work with a continuous-action method is needed; the six-run SAC benchmark remains unreplicated.

\section{Public fixed-condition Parkinsonian tremor traces, not adaptive policy data}
PhysioNet hosts human Parkinsonian finger-tremor velocity recordings from 16 people implanted with DBS, under effective DBS on/off and dopaminergic medication on/off conditions: \url{https://physionet.org/content/tremordb/1.0.0/}. We fetched the 55 original plain-text recordings in these eight resting-condition categories and numerically analyzed 15 represented subjects. Header file rofh.hea lists v5rof.rit, while the actual directory and source description list v5rof.let; the script explicitly corrects that one source-header discrepancy instead of silently losing a trace. Welch power in the 4--6 Hz tremor band was calculated on each ~60-second, 100-Hz raw instrument trace and paired within subject at a fixed medication condition, yielding 25 DBS-on/DBS-off pairs. The descriptive median on/off power ratio across all pairs was 0.873; stratifying by medication gave 0.805 for 13 medication-on pairs and 1.332 for 12 medication-off pairs. These simple ratios are heterogeneous and conflict with any uniform suppression claim. Instrument ranges and calibration differ by subject/recording, so the paired ratio is more interpretable than absolute pooled power but still not causal: recordings may differ in order, behavior, recording quality, and medication state is not randomized here. Repeated recordings are not new patients. Most importantly, this is finger velocity, not DBS-Gym low-beta LFP; stimulation is fixed per condition, not a time-indexed adaptive action. It cannot supply logged state-action-next-state trajectories or validate a reinforcement policy. Source recording hashes, subject-condition mapping and all individual ratios are in results/parkinson\_tremordb.json; code src/parkinson\_tremordb.py. This genuine human-source audit adds one Parkinson-specific scientific data product to the strict tool ledger (now nine), not an action-policy benchmark.

\section{What the published clinical offline-RL experiment measures, and what this project lacks}
Gao and colleagues' 2023 clinical offline-RL study is a useful boundary case for the present exercise, rather than a result we can claim to reproduce. Their four participants had implanted sensing-and-stimulation devices and monthly clinical visits. According to their report, a logged experience buffer preserved the LFP signal and stimulation amplitude over time; efficacy checks also used bradykinesia tasks, patient ratings and wearable tremor. Those streams correspond to different temporal resolutions. LFP and action are available at repeated control steps, whereas the clinical checks can be sparse or end-of-session. The authors describe amplitude as a fraction of clinician-set conventional DBS, with zero and one as endpoints, and use offline policy training, compressed policies and model-based off-policy evaluation to rank candidates before clinical testing. Their stated outcome is comparable symptom control with lower stimulation energy in their four-person investigation. Their clinical report is \url{https://arxiv.org/html/2302.02477v4}; these are its findings, not an outcome of our models.

Our 55 PhysioNet tremor recordings do not have that logged joint structure. Each recording belongs to a fixed on/off stimulation and medication condition, and its velocity trace runs through time, but it does not provide the contemporaneous changing stimulation amplitudes, implant LFP, physician-selected amplitude, trial-level feedback or the next state under a dynamically chosen action. An on/off recording pair is not a sequential intervention history: computing a ratio after recording does not reconstruct which action a hypothetical adaptive policy would have taken at any sample. Likewise, our GEO blood sample accessions label disease status rather than within-patient stimulation effects. The same person can contribute several condition recordings, so neither the 55 signals nor the 25 matched ratios are 55 or 25 independent adaptive-control evaluations. This gap remains even though the sample-count gate exceeds 120.

The DBS-Gym paper supplies a different kind of evidence. It specifies a configurable oscillator network, with Env0 emphasizing beta-band behavior, Env1 adding spatial complications and Env2 adding drift, and reports multi-evaluation controller results under these simulator choices: \url{https://arxiv.org/html/2505.09624}. Our exact upstream Env0 OFF and high fixed controls reached the full horizon at one seed. Recomputing the authors' filtered spectral statistic on those cached signals gave a 27.536\% high/OFF ratio, while their Table 1 reports 19.8\% with standard deviation 1.9 across six evaluations. This difference cannot be explained by substituting our 22.48\% mean-window reward-component ratio, because its aggregation differs from the published filtered spectrum. Nor is this discrepancy a clinical failure, since both are simulator outputs. The published benchmark comparison needs its evaluation seed/configuration protocol and a trained controller run through the same statistic, ideally with paired uncertainty. Our short neural policy instead collapsed to OFF, so extending that policy to a 5555-step rollout would consume computation without curing its demonstrated failure. A serious next within-project pivot would change training and evaluation protocol first, then run matched, multi-seed full-length comparisons, reporting all attempted policies rather than selecting the best seed.

This evidence hierarchy also clarifies what clinical validation cannot mean here. An independent blood-diagnosis cohort can stress-test a molecular classifier, not a stimulation controller. A fixed-condition human finger trace can describe heterogeneity of on/off tremor power, not counterfactual adaptive actions. A native simulator benchmark can compare algorithms under specified dynamics, not prove safety for implant programming. Only properly governed action-linked patient trajectories and clinical assessment can begin to bridge the final gap; this project has no access to such trajectories. Until then all stimulation values in this report are engineering controls in simulation, not patient dose recommendations.

\section{Sources}
\begin{enumerate}
\item Kuzmina E. et al., KDD 2025, DBS-Gym code. \url{https://github.com/NevVerVer/DBS-Gym}.
\item Publisher record: \url{https://dl.acm.org/doi/10.1145/3711896.3737397}.
\item GSE99039 GEO accession page: \url{https://www.ncbi.nlm.nih.gov/geo/query/acc.cgi?acc=GSE99039}.
\item Full checksum and accession-level manifest: data/parkinson/provenance.json and data/parkinson/accession\_features.csv.
\end{enumerate}
\end{document}
